% !TEX root = ../main.tex
\chapter{General Presentation and Context of the Internship}
\label{chap:presentation}

\section*{Introduction}
\phantomsection
\addcontentsline{toc}{section}{Introduction}

The purpose of this chapter is to situate the project in its general framework, and
to set out the problem it addresses, the objectives it pursued, and the manner in
which the work was conducted.

The work reported here was carried out as a research internship at the AI for
Decision Making Lab of Concordia University, under the supervision of
Dr.~Yassine Yaakoubi, and constitutes the research component of the final year of
engineering studies at the Ecole Polytechnique de Tunisie.()


We first present the host organisation, its structure and its areas of research.
We then state the problem the work addresses, in the terms in which it arises for
a deployed system, and situate it in the systems already deployed and the
regulations that govern them, before setting out the objectives pursued and the
deliverables produced. The chapter closes with the methodology by which the work
was conducted.

\section{Host organisation}
\label{sec:ch1-host-organization}

The internship was hosted by the \textbf{AI for Decision Making Lab (AI4DM)}, a
research group attached to the Gina Cody School of Engineering and Computer
Science at Concordia University, in Montreal, Canada. Table~\ref{tab:host-organization}
summarises the host organisation; the subsections below present the host
university, the faculty to which the laboratory belongs, and finally the research
group in which the work was carried out.

\begin{table}[H]
\centering
\caption{Host organisation details}
\label{tab:host-organization}
\begin{tabular}{|l|>{\raggedright\arraybackslash}m{9cm}|}
\hline
\textbf{Institution}   & Concordia University \\ \hline
\textbf{Faculty}       & Gina Cody School of Engineering and Computer Science \\ \hline
\textbf{Department}    & Mechanical, Industrial and Aerospace Engineering (MIAE) \\ \hline
\textbf{Laboratory}    & AI for Decision Making Lab (AI4DM) \\ \hline
\textbf{Supervisor}    & Dr. Yassine Yaakoubi \\ \hline
\textbf{Location}      & Montreal, Quebec, Canada \\ \hline
\textbf{Research axes} & Optimisation, learning and simulation for large-scale
                         decision making under uncertainty; large language models
                         and retrieval-augmented generation; uncertainty
                         quantification and reliability of AI systems \\ \hline
\end{tabular}
\end{table}

\subsection{Concordia University}
\label{subsec:ch1-concordia}

Concordia University is a public comprehensive research university located in
Montreal, Quebec. It was founded in 1974 by the merger of two older
institutions\citeN{net-merger}: Loyola College, established in 1896, and Sir
George Williams University, which grew out of the evening classes offered by the
Montreal YMCA from 1851\citeN{net-sgw}. Its name comes from the motto of the
city of Montreal, \emph{Concordia salus}, that is to say ``well-being through
harmony''.

With 43,922 students enrolled in credit courses in 2024--2025 over its two
campuses, Sir George Williams downtown and Loyola in the west of the city,
Concordia is today one of the largest universities in Canada and one of the
leading Canadian institutions in engineering and applied
sciences\citeN{net-fastfacts}.
% NOTE (figure changed, 2026-08-13): this sentence previously read 48,657.
% That number is the university's GRAND TOTAL and includes 4,735 non-credit
% Continuing Education registrations; the sentence describes degree students,
% so the credit-enrolment figure (43,922) is the correct one. Both appear on
% the same official page (net-fastfacts): credit 43,922, grand total 48,657.

\subsection{Gina Cody School of Engineering and Computer Science}
\label{subsec:ch1-gina-cody}

The Gina Cody School of Engineering and Computer Science is the faculty of
Concordia University dedicated to engineering and computer science. In 2018, it
became the first engineering faculty in Canada to be named after a woman,
following a donation of fifteen million dollars from Gina Parvaneh
Cody\citeN{net-codygift}, who had been in 1989 the first woman to obtain a
doctorate in building engineering at Concordia, and who has been the chancellor
of the university since January 2025\citeN{net-codyprofile}.

The school gathers 11,315 students and 262 faculty members around 43
undergraduate and graduate programmes and 20 interdisciplinary research
centres\citeN{net-ginacody}.
% NOTE (discrepancy between two official Concordia pages, checked 2026-08-13):
% the figure 11,315 is the Gina Cody School's OWN "At a glance" page
% (net-ginacody), which is what is cited here. The university-level Fast Facts
% page (net-fastfacts) separately reports Gina Cody School credit enrolment as
% 11,576 for 2024-25. The two are not reconciled by either page. If questioned,
% the answer is that the school's own published figure is cited.
It is organised into seven departments, among which the Department of Mechanical,
Industrial and Aerospace Engineering (MIAE), to which the host laboratory belongs.
Its teaching and research activities are housed principally in the Engineering,
Computer Science and Visual Arts Complex shown in Figure~\ref{fig:ev-building}.

\begin{figure}[H]
\centering
\includegraphics[width=0.7\textwidth]{ev-building.jpg}
\caption{The EV Complex, home of the Gina Cody School (photograph: Tasitch, CC BY 2.0)}
\label{fig:ev-building}
\end{figure}

\subsection{Research group, supervision, and research axes}
\label{subsec:ch1-research-group}

The work presented in this report was conducted within the AI for Decision Making
Lab (AI4DM), the research group led by Dr.~Yassine Yaakoubi, assistant
professor in the MIAE department.

The research of the group is situated at the intersection of optimisation,
learning and simulation for large-scale real-world decision making. It is
organised around two main axes: machine-learning-augmented optimisation, which
integrates learned models into classical optimisation methods, and end-to-end
optimisation learning. To this end, the group combines mixed-integer and
stochastic programming, decomposition methods, deep learning, reinforcement
learning and graph neural networks, in order to design decision-support systems
that are efficient, robust and scalable, with applications in logistics,
transportation and sustainable supply chains.

A further strand of the group's activity extends this concern with robustness to
systems built on large language models. It covers retrieval-augmented generation,
where a model is grounded in evidence drawn from an external corpus, and the
uncertainty quantification and reliability of such systems, that is, the question
of how far their outputs can be trusted and how that trust can be measured. The
work reported here belongs to this strand.

This scientific environment, centred on decision making under uncertainty and on
the robustness of the resulting systems, is directly aligned with the subject of
the internship, which aims at providing reliability guarantees for multi-agent
question-answering systems.

\section{Problem statement}
\label{sec:ch1-problem-statement}




\section{Objectives and expected deliverables}
\label{sec:ch1-objectives}


\section{Working methodology}
\label{sec:ch1-methodology}


\section*{Conclusion}
\phantomsection
\addcontentsline{toc}{section}{Conclusion}

